\section{Conclusion}

Recent AI milestones, from solving Millennium Prize Problems to committing cybercrimes, suggest that the gravest risks of superintelligence will result from models' collective behaviour as they expand and accelerate through digital spaces. Calls to `pace the frontier' and slow AI progress also reveal the research community's desperate need for tools to ensure this collective intelligence is safe and aligned with human values. As experts in the social dynamics and group behaviour of both humans and AI agents, Artificial Societies aims to help researchers address these risks by applying insights from a century of social science to the frontier of machine learning.

This blog post takes an initial step in that direction. We distinguish individual alignment from group alignment, and define the latter in a way that clarifies how ACI can support human values and intentions even when humans cannot guarantee the specification of each agent. Our numerical experiments expand on this definition to provide a concrete mechanism for achieving group alignment via `mixed swarms' of frontier models and agentic models of human behaviour. 

The ABMs' results communicate two potentials. First, mixed swarms can address fundamental concerns about propagating human information and ensuring our values aren't crowded out by fast-moving agent swarms, so long as human society also changes how we interact and share content from AI. Second, if we fail to take appropriate safety precautions around agent governance and interaction in shared digital environments, these models may rapidly dominate the information space and make it impossible for humans to coexist alongside them. On the internet, the socioeconomic implications would be alarming, so safety researchers and AI leaders must rigorously consider the risks of misaligned collective intelligence. 

As next steps, we are scaling up our numerical ABMs to true multi-agent RL environments where we can deploy open-weight frontier models alongside our foundational models of human behaviour, observe their interactions, and apply group-level reward functions to the frontier models. We are also committing to publishing and open-sourcing all our research on aligned collective intelligence, since AS should not be the only team doing work in this consequential domain. In this light, we encourage independent researchers, labs, and our fellow industry builders to explore and collaborate on advancing group alignment. Without more joint efforts, ACI may irreversibly damage modern society. 